\section{Actual local line coefficients and intrinsic curve degree}

\begin{definition}[The coefficient in an actual original affine frame]
\label{def:fc-original-line-section-frame-coefficient}
\leanok
\lean{EnvelopeProperty.CurveLineSectionOrders.frameCoefficient, EnvelopeProperty.CurveLineSectionOrders.frameCoefficient_ne_zero}
\uses{thm:fc-original-line-zero-ideal-module}
For the original section $s:O_X\to L$ and an actual affine frame
$e:L|_U\simeq O_U$, let $c_e(s)$ be the original section of $O_X(U)$
given by applying $e$ to $s(1)$. If $X$ is integral and $s\ne0$, then
$c_e(s)\ne0$ on every nonempty original affine framed open, and the
actual original zero ideal on $U$ is exactly $(c_e(s))$.
\end{definition}
\begin{proof}
The whole image defining the original zero ideal is the coefficient
principal ideal in an actual frame. The nonzero original section
internally makes that original zero ideal Cartier. Its original stalk
ideal at any point of a nonempty framed open contains a nonzero
Cartier equation. If the original coefficient were zero, the actual
localization of its principal ideal would be zero, a contradiction.
\end{proof}

\begin{theorem}[The actual frame coefficient has the original local order]
\label{thm:fc-original-line-section-frame-order}
\leanok
\lean{EnvelopeProperty.CurveLineSectionOrders.localOrder_eq_frameCoefficient_order, EnvelopeProperty.CurveLineSectionOrders.weightedLocalOrder_eq_frameCoefficient_order}
\uses{def:fc-original-line-section-frame-coefficient, thm:fc-curve-cartier-weighted-degree}
On an original integral locally Noetherian scheme, the native local
order of the same original zero ideal at $x\in U$ is exactly
$\operatorname{ord}_x(c_e(s))$ under the original generic germ. Its
weight over the original field is this same integer times the degree
of the actual original residue field.
\end{theorem}
\begin{proof}
The preceding actual coefficient is a nonzero generator of the same
original Cartier ideal on the actual affine open. The genuine
local-equation order comparison identifies its native scheme order
with the zero ideal's native order. Multiplication by the literal
original residue-field degree gives the weighted formula.
\end{proof}

\begin{theorem}[True section orders and generated line degrees are nonnegative]
\label{thm:fc-original-curve-generated-line-degree-nonnegative}
\leanok
\lean{EnvelopeProperty.CurveLineSectionOrders.localOrder_nonneg, EnvelopeProperty.CurveLineSectionOrders.degree_nonneg_of_nonzero_section, EnvelopeProperty.CurveLineSectionOrders.degree_nonneg_of_generatingSections, EnvelopeProperty.CurveLineSectionOrders.degree_pullback_nonneg_of_generatingSections}
\uses{thm:fc-original-line-section-frame-order, thm:fc-original-curve-line-degree-zero-section}
On an original integral locally Noetherian scheme every original
Cartier ideal's native local order is nonnegative. Consequently on
any original integral projective curve, the intrinsic degree of an
original line with an actual nonzero section is nonnegative. In
particular this holds for any actual globally generated line and,
after any literal curve map, for the actual pullback of any original
globally generated line. No normality, smoothness, prescribed degree,
nefness, or section-order certificate is assumed.
\end{theorem}
\begin{proof}
At an actual codimension-one point, the native order of a nonzero
regular Cartier equation is the nonnegative length of the same
principal quotient of the actual original stalk. At all other points
its native scheme order is zero. The original residue-field degree is
nonnegative. The genuine full Euler degree formula for the original
nonzero section is therefore a finite-support sum of nonnegative
weighted orders. An actual generating family contains a nonzero
original section because an invertible line on the nonempty original
scheme is nonzero. Original pullback preserves the actual generating
family, so the same argument applies to any actual curve pullback.
\end{proof}
